\chapter*{\centering KATA PENGANTAR}
\addcontentsline{toc}{chapter}{KATA PENGANTAR}

\setstretch{1.15}  % Mengurangi spasi dalam bagian ini saja
\setlength{\parskip}{0.3em} % Mengatur jarak antar paragraf

Puji syukur selalu terpanjat kehadirat Allah SWT atas segala rahmat, petunjuk, dan kesempatan yang diberikan sehingga penulis dapat menyelesaikan proposal tugas akhir dengan judul ``Rancang Bangun Instrumentasi Sensor \textit{Tunneling Magnetoresistance} Berbasis Nanofiber \ch{Fe3O4}/PVA-Sitrat-ADH untuk Deteksi Formalin pada Bakso Menggunakan Komparasi Model Klasik (SVM, Random Forest) dan Kuantum (QSVC, VQC)''.

Penulis menghaturkan terima kasih dan apresiasi kepada seluruh pihak yang telah berkontribusi positif terhadap penulis. Ucapan terima kasih penulis haturkan kepada:

\begin{enumerate}[noitemsep, topsep=2pt, parsep=0pt, partopsep=0pt]
    \item Allah Subhanahu Wa Ta'ala, karena atas nikmat, karunia, dan izin-Nya lah, Alhamdulillah proposal tugas akhir ini selesai disusun.
    \item Ibu Ara Kusmirah serta segenap keluarga yang tidak henti memberikan dukungan secara moril dan material.
    \item Bapak Mada Sanjaya W.S., M.Si., Ph.D., selaku Dosen Pembimbing Tugas Akhir yang senantiasa memberikan ilmu, arahan, bimbingan, hingga motivasi selama proses penelitian dan penulisan proposal tugas akhir ini.
    \item Segenap civitas akademika Fisika UIN SGD Bandung yang telah menyisihkan waktunya dalam memberikan pengetahuan, diskursus, dan motivasi yang berarti bagi penulis, khususnya dalam mempelajari keilmuan dalam menyelesaikan proposal tugas akhir.
    \item Alma Triana Nabilla, yang tidak henti membersamai, memberikan dukungan penuh secara moril dan material setiap waktunya, sehingga penulis memiliki motivasi tinggi dalam menyelesaikan studi strata I. Semoga hal-hal baik senantiasa beriring secara simultan di antara ruang dan waktu padanya.
    % TODO(Fahry): tambahkan di sini pihak lain yang ingin diucapkan terima kasih
    % (dosen pembimbing akademik, rekan penelitian/lab, sahabat, dsb.) sesuai
    % dengan orang-orang yang sebenarnya terlibat dalam proses studimu.
\end{enumerate}

Dengan penuh kerendahan hati, penulis menyadari bahwa proposal ini turut jauh dari kata sempurna dan memiliki banyak keterbatasan. Karenanya, penulis sangat mengharapkan kritik serta saran yang membangun. Semoga proposal ini dapat memberikan manfaat bagi penulis pribadi, maupun bagi pengembangan ilmu pengetahuan, teknologi, serta pihak-pihak yang memerlukannya.

\vspace{0.2cm}
\begin{flushright}
    Bandung, 23 September \the\year{}\\[1.2cm]
    \textbf{Fahry Rizky Samsudin}
\end{flushright}
